\documentclass[11pt, a4paper]{article}
\usepackage[a4paper, top=2.5cm, bottom=2.5cm, left=2cm, right=2cm]{geometry}
\usepackage{fontspec}
\usepackage[english, bidi=basic, provide=*]{babel}
\babelprovide[import, onchar=ids fonts]{english}
\babelfont{rm}{TeX Gyre Termes}
\babelfont{sf}{TeX Gyre Heros}
\babelfont{tt}{TeX Gyre Cursor}
\usepackage{enumitem}
\usepackage{amsmath}

\title{Study Notes: Meiosis and Sexual Life Cycles (Chapter 13)}
\author{Subject: Biology}
\date{}

\begin{document}
\maketitle

\section*{Core Concepts}
\begin{itemize}
    \item \textbf{Heredity:} The transmission of traits from parents to offspring.
    \item \textbf{Variation:} Demonstrated by the differences in appearance that offspring show from parents and siblings.
    \item \textbf{The Link:} Offspring acquire genes from parents by inheriting chromosomes.
    \item \textbf{Asexual vs. Sexual Reproduction:}
    \begin{itemize}
        \item \textbf{Asexual:} Single parent; offspring are \textbf{clones} (genetically identical).
        \item \textbf{Sexual:} Two parents; offspring have unique combinations of genes; genetic variation is high.
    \end{itemize}
\end{itemize}

\section*{Key Definitions}
\begin{itemize}
    \item \textbf{Gene:} Hereditary unit made of DNA.
    \item \textbf{Locus:} A gene's specific location on a chromosome.
    \item \textbf{Gametes:} Reproductive cells (sperm and eggs) that transmit genes.
    \item \textbf{Somatic Cells:} All cells of the body except gametes (diploid, $2n = 46$ in humans).
    \item \textbf{Karyotype:} Ordered display of pairs of chromosomes from a cell.
    \item \textbf{Homologous Chromosomes (Homologs):} Two chromosomes of a pair with the same length, centromere position, and genes controlling the same characters.
    \item \textbf{Autosomes:} Chromosomes 1--22 in humans (non-sex chromosomes).
    \item \textbf{Sex Chromosomes:} X and Y; determine the individual's sex.
\end{itemize}

\section*{Important Formulas / Relations}
\begin{itemize}
    \item \textbf{Diploid Number ($2n$):} For humans, $2n = 46$.
    \item \textbf{Haploid Number ($n$):} For humans, $n = 23$.
    \item \textbf{Independent Assortment Combinations:} $2^n$. For humans, $2^{23} \approx 8.4$ million.
    \item \textbf{Random Fertilization Combinations:} $2^n \times 2^n$. For humans, $8.4 \text{ million} \times 8.4 \text{ million} \approx 70$ trillion.
\end{itemize}

\section*{Processes / Workflows}

\subsection*{The Stages of Meiosis}
Meiosis is preceded by DNA replication in \textbf{Interphase}. It consists of two consecutive divisions:
\begin{enumerate}
    \item \textbf{Meiosis I (Reductional Division):}
    \begin{itemize}
        \item \textbf{Prophase I:} Homologs pair up (\textbf{synapsis}) and exchange segments (\textbf{crossing over}). \textbf{Chiasmata} form.
        \item \textbf{Metaphase I:} Homologous pairs align at the metaphase plate; \textbf{Independent Assortment} occurs.
        \item \textbf{Anaphase I:} Homologs separate; sister chromatids remain attached at centromeres.
        \item \textbf{Telophase I \& Cytokinesis:} Two haploid cells form.
    \end{itemize}
    \item \textbf{Meiosis II (Equational Division):} Similar to mitosis; sister chromatids separate.
    \begin{itemize}
        \item Results in \textbf{four haploid daughter cells}, each genetically distinct.
    \end{itemize}
\end{enumerate}

\subsection*{Variety of Sexual Life Cycles}
\begin{itemize}
    \item \textbf{Animals:} Gametes are the only haploid cells. Zygote (2n) divides by mitosis.
    \item \textbf{Plants (Alternation of Generations):} Includes both diploid multicellular (\textbf{sporophyte}) and haploid multicellular (\textbf{gametophyte}) stages.
    \item \textbf{Fungi/Protists:} Zygote is the only diploid stage; it immediately undergoes meiosis to produce haploid cells.
\end{itemize}

\section*{Common Mistakes and Clarifications}
\begin{itemize}
    \item \textbf{Sister Chromatids vs. Homologs:} Sister chromatids are identical copies of one chromosome; homologs are a pair of similar chromosomes (one from each parent).
    \item \textbf{Replication timing:} DNA is replicated \textit{once} before Meiosis I, but \textit{never} between Meiosis I and II.
    \item \textbf{Reduction vs. Equation:} Meiosis I reduces the \textit{chromosome set count} (2n to n); Meiosis II reduces the \textit{chromatid count} but keeps the set count at n.
\end{itemize}

\section*{Exam-Focused Quick Revision}
\begin{itemize}
    \item \textbf{Three Sources of Genetic Variation:}
    \begin{enumerate}
        \item \textbf{Independent Assortment:} Random orientation of homologs at Metaphase I.
        \item \textbf{Crossing Over:} Exchange of DNA between nonsister chromatids in Prophase I, creating \textbf{recombinant chromosomes}.
        \item \textbf{Random Fertilization:} Any sperm can fuse with any egg.
    \end{enumerate}
    \item \textbf{Synapsis Mechanism:} The \textbf{synaptonemal complex} (zipper) holds homologs together.
    \item \textbf{Evolutionary Significance:} Variation is the raw material for natural selection.
\end{itemize}

\end{document}